\section{Limitations and Future Work}
\label{sec:limits}

\paragraph{Evidence} Only the completion gate and the router are measured here, and the gate only on constructed end states. End-to-end rates and latencies with local models remain to be measured (Section~\ref{sec:setup}); latency will depend strongly on the host and on which models are installed.

\paragraph{Rule-based verification} $\Phi$ is keyed to request wording and a few search engines, and Section~\ref{sec:results} lists its blind spots. We plan to derive typed step predicates from the planner's expected outcomes, add a local vision-model check of the final screenshot as a second opinion, and calibrate both on labelled trajectories.

\paragraph{Small models and hand-written rules} Planning with 1.5B--7B models is weak, and deterministic overrides compensate on common patterns, some of them site-specific. The extraction handler also contains post-processing written for one case-study site. \todo{AUTHOR: remove or generalize the case-study-specific extraction code in \texttt{backend/agent/nodes.py} (lines 1108--1139) before any end-to-end evaluation, or disclose it here.} Canvas- and WebGL-rendered pages expose no useful DOM and depend on the vision path.

\paragraph{Inactive mechanisms and roadmap} Strategy-specific recovery, memory retrieval, the routing cool-down and the verification switch need small fixes to take effect (Table~\ref{tab:status}); the ablations are meant to measure them once they do. For the local-system extension we will implement the path and process policy, the predicate library, an AT-SPI provider (testable headlessly on Linux) and the \texttt{action\_gate} node, then Windows UI Automation. Longer term, we plan to split planning, grounding and verification across cooperating agents, with the verifier given no write access to the environment.
